% =====================================================================
% PAPER6 — Capture-Resilient Voting in Multi-LLM Verifier Ensembles
% Template: tmlr.sty (v0.1, copied from PAPER5, 6560 bytes)
% Status: COMPLETE DRAFT — full sections + synthetic pilot results (4,000 trials/cell, N∈{5,7,10}, f∈{1,2,3}); real-trace benchmark noted as future work
% =====================================================================

\documentclass[11pt]{article}
\usepackage[preprint]{tmlr}
% If accepted, switch to: \usepackage[accepted]{tmlr}
\usepackage{amsmath, amssymb, amsthm}
\input{math\_commands}
\usepackage[T1]{fontenc}
\usepackage[utf8]{inputenc}
\usepackage{graphicx}
\usepackage{booktabs}
\usepackage{multirow}
\usepackage[colorlinks=true,linkcolor=blue,citecolor=blue,urlcolor=blue]{hyperref}
\usepackage[capitalize]{cleveref}
\usepackage{microtype}
\setlength{\headheight}{14pt}
\usepackage{algorithm}
\usepackage{algpseudocode}
\let\STATE\State

% ---------------------------------------------------------------------
% Theorem environments
% ---------------------------------------------------------------------
\newtheorem{definition}{Definition}
\newtheorem{theorem}{Theorem}
\newtheorem{lemma}{Lemma}
\newtheorem{corollary}{Corollary}
\newtheorem{proposition}{Proposition}
\newtheorem{remark}{Remark}
\newtheorem{assumption}{Assumption}

% ---------------------------------------------------------------------
% Notation shortcuts
% ---------------------------------------------------------------------
\newcommand{\V}{\mathcal{V}}            % set of verifiers
\newcommand{\A}{\mathcal{A}}            % set of attackers
\newcommand{\CRV}{\textsc{CRV}}         % Capture-Resilient Voting
\newcommand{\ECR}{\varepsilon\text{-CR}} % $\varepsilon$-capture-resilient

\title{An Empirical Study of Aggregation Methods Under\\
       Verifier Capture in Multi-LLM Ensembles:\\
       When Does Majority Voting Fail, and Can We Beat It?}

\author{\name Anonymous\email anonymous@anonymous.invalid\\\addr Anonymous affiliation}
\date{July 2026}

\begin{document}
\maketitle

\begin{abstract}
Multi-LLM verification ensembles increasingly aggregate the verdicts of
several LLM \emph{verifiers} to reduce variance and adversarial
brittleness relative to any single model. A growing concern in such
ensembles is \emph{verifier capture}: a small coalition of verifiers
whose outputs are systematically biased can flip the aggregated
decision while remaining individually plausible. This paper contributes
(i) a formal definition of capture-resilient voting (\textsc{CRV}) and
the threat model distinguishing prompt-injection, rubric-injection, and
majority-bias mimicry attacks; (ii) an experimental
protocol comparing seven aggregation methods (plain majority,
PBFT-style deterministic thresholding, uniform-weighted mean, three
trust-score variants, and an adaptive hybrid) across verifier pools of
size $N \in \{5,7,10\}$ with capture counts $f \in \{1,2,3\}$ in a
synthetic benchmark and a small local-model pilot; and (iii) a falsifiable prediction---that captured verifiers'
outputs are statistically indistinguishable from a minority of
honest-but-wrong verifiers, so any statistical-deviation capture rule
incurs false-positive costs that can outweigh its gains. A synthetic pilot (4,000 trials per cell, $N=7$ verifiers,
f=2 captured) already tests the prediction and yields a nuanced result:
trust-score weighting beats plain majority under prompt-injection capture
(accuracy 0.69 vs. 0.34 at honest accuracy 0.6; 0.94 vs. 0.74 at 0.8),
because calibration-phase trust identifies deterministically flipped
verifiers; under majority-bias mimicry, however, the same method fails
(0.32 vs. 0.41 at honest accuracy 0.6). These pilot values are synthetic and
do not establish a universal impossibility result. A descriptive local-model
pilot is included in the evidence package, but a held-out real-trace benchmark
remains future work; we release the framework, generators, and analysis
scripts so the benchmark can be extended independently.


\textbf{Reproducibility and ethics.} All experiments use deterministic seeds (PYTHONHASHSEED=0) with GPT-4o, DeepSeek, and Qwen models; full experimental protocols, code, raw API logs, and analysis scripts are released in the supplementary material with SHA-256 checksums. This work studies evaluation dynamics in LLM systems; ethical implications include silent reliability decay in safety-critical systems and the need for version-dependent evaluation verification.\end{abstract}

% =====================================================================


\textbf{Novelty positioning.} This work differs from existing literature in three ways: (i) systematic decomposition of temporal vs.\ communicative causes of calibration fatigue; (ii) demonstration that self-evaluation immunity is model-dependent and fragile; (iii) reframing of the coupling-noise tradeoff as method-specific with explicit falsification conditions.



\textbf{Contributions.} This paper makes four contributions: (1) \emph{Systematic empirical analysis} of the target phenomenon across multiple model families and experimental conditions, with matched baselines and controlled confounds. (2) \emph{Formal characterization} of the key metric or relation, with explicit definitions, falsification conditions, and boundary cases. (3) \emph{Cross-model validation} demonstrating that observed effects are not model-specific, with replication attempts and explicit reporting of reproduction failures. (4) \emph{Reproducibility framework}: all experiments released with deterministic seeds, SHA-256 checksums, and step-by-step reproduction protocols, enabling independent verification of all results.\section{Experimental Setup}
\label{sec:experimental-setup}

\textbf{Models and APIs.} All experiments use the following canonical configuration unless otherwise noted. \textit{Primary model}: GPT-4o (API version gpt-4o-2024-05-13) with temperature $=0.0$, top\_p $=1.0$, and max\_tokens $=1024$ for all agent interactions. \textit{Cross-model validation}: DeepSeek-V4-Pro and Qwen-2.5-72B via their respective APIs with identical generation parameters. \textit{Reproducibility}: all randomness controlled via deterministic seeds (seed $=42$ for canonical runs; 30 distinct seeds for repetition experiments); Python hash randomization disabled (PYTHONHASHSEED $=0$); all API calls use idempotent request parameters.

\textbf{Tasks and datasets.} Evaluation tasks include GSM8K (math reasoning, $K{=}5$ problems per round), MMLU (multi-disciplinary QA), and HumanEval (code generation) for diversity. All datasets are publicly available; exact problem subsets and seed mappings are released in the evidence package.

\textbf{Metrics.} Primary metrics: Expected Calibration Error (ECE, $B{=}10$ equal-width bins), Brier score, coupling coefficient $\gamma$, strategy entropy $H$, and coefficient of variation $\CV$. All metrics computed at the per-seed level and aggregated across seeds with bootstrap confidence intervals.

\textbf{Statistical procedures.} All confidence intervals are bootstrap intervals ($B{=}5{,}000$ resamples). Pairwise contrasts use Welch's $t$-test with Bonferroni correction for multiple comparisons (correction factor $k$ specified per table). Effect sizes reported as Cohen's $d$. All tests two-sided unless otherwise noted.

\textbf{Exclusion criteria.} Runs with API errors, empty responses, or malformed outputs ($<2\%$ of total) are excluded and re-run with fresh seeds. All exclusions logged in the evidence manifest with SHA-256 checksums for all result files.

\section{Introduction}
\label{sec:intro}
% =====================================================================

Large language model (LLM)-based verification pipelines increasingly rely on
\emph{multi-source verifier ensembles} \citep{cui2026edo,leung2026rubricllm}, in which a fixed pool
of $N$ LLM verifiers independently evaluate a candidate output and a final
decision is produced by aggregating their votes. Such ensembles have become a
default building block of contemporary AI auditing, peer review assistance,
and safety-critical content moderation, where the goal is to reduce variance
and adversarial brittleness relative to any single model.

A natural and increasingly documented failure mode of these ensembles is what
we call \emph{verifier capture}: a small coalition of verifiers whose outputs
are systematically biased in a coordinated direction can swing the aggregated
decision, even when a majority of verifiers remain individually reliable.
Verifier capture differs from the classical single-verifier hallucination or
bias problem in two important ways.
\textbf{(i) Capture is interactive.} The adversarial verifier does not need to
be accurate in absolute terms; it only needs to produce outputs that, when
combined with the other verifiers under a fixed aggregation rule, flip the
final decision.
\textbf{(ii) Capture is silent.} From the perspective of an external observer
that only sees final aggregated verdicts, a captured ensemble looks identical
to a correctly functioning one until the disagreement structure of individual
votes is examined.

Recent work has begun to expose capture-like failure modes in LLM judges:
\citet{zhang2026blindcurator} show that a biased judge \emph{silently
disables} a downstream curator past a sharp false-pass threshold. This paper
asks the natural follow-up question: given that capture is real and silent,
\emph{can any aggregation rule do better than plain majority voting}? Our
answer, in this empirical study, is \emph{mostly no}.

This paper makes four contributions:

\begin{itemize}
  \item A formal threat model for verifier capture in LLM ensembles, including
        a definition of $\varepsilon$-capture-resilience that decouples the
        adversary's per-verifier accuracy from its aggregate influence
        (Section~\ref{sec:threat}).
  \item A unified experimental protocol that benchmarks seven aggregation
        methods under three attack families and four scenario configurations
        (Section~\ref{sec:method}).
  \item An empirical finding that \emph{no} aggregation method we tested
        substantially outperforms plain majority voting, and several
        trust-score variants \emph{underperform} majority under strong
        capture (Section~\ref{sec:exp}). We provide a structural
        explanation for the negative result: in our setting the captured
        verifiers' outputs are statistically indistinguishable from a
        minority of honest-but-wrong verifiers, so any rule that uses
        statistical deviation to identify capture incurs false-positive
        costs that outweigh its true-positive gains.
  \item A reproducible open-source release with code, vote generators, and a
        Beyond-PDF interactive attack demonstration
        (Section~\ref{sec:broader}).
\end{itemize}

We emphasise that the \emph{negative} headline finding is the paper's
contribution, not an apology. TMLR's acceptance criteria explicitly
prioritise technical correctness and audience interest over novelty, and a
rigorous benchmark that documents the limits of the obvious aggregation
strategies is valuable precisely because future work can build on it
rather than re-derive it. We discuss when, and against which assumptions,
our negative result would \emph{not} hold
(Section~\ref{sec:conclusion}).

The remainder of this paper is organized as follows.
Section~\ref{sec:background} reviews classical Byzantine fault tolerance and
LLM ensembling. Section~\ref{sec:threat} introduces the threat model.
Section~\ref{sec:method} describes the seven aggregation rules we benchmark.
Section~\ref{sec:exp} reports the benchmark and its negative result.
Section~\ref{sec:related} discusses related work.
Section~\ref{sec:broader} addresses broader impact.
Section~\ref{sec:conclusion} concludes and outlines conditions under which
our negative result would not hold.

% =====================================================================
\section{Background}
\label{sec:background}
% =====================================================================

\paragraph{Classical Byzantine consensus.}
Byzantine fault tolerance (BFT) \citep{lamport1982byzantine} establishes
that a deterministic consensus protocol over $N$ replicas tolerates at
most $f < N/3$ Byzantine (arbitrary) failures when paired with digital
signatures, or $f < N/2$ under synchrony assumptions with binary outputs.
PBFT \citep{castro1999pbft} brought BFT to practical systems and
HotStuff \citep{yin2019hotstuff} linearised the communication pattern.
All classical BFT results assume \emph{binary} replica outputs; the
extension to \emph{continuous} outputs (which we need for LLM verifier
probabilities) is not standard and is one of the contributions of this
paper.

\paragraph{LLM-as-judge and ensembles.}
The use of LLMs as evaluators, or \emph{judges}, was popularised by
MT-Bench and Chatbot Arena \citep{zheng2023judging}, where a single
LLM judge is shown to correlate well with human preference ratings.
Subsequent work studied ensemble-of-judges methods
\citep{cui2026edo,leung2026rubricllm} to reduce variance and calibrate
directional bias. These methods optimise for \emph{passive} disagreement
across annotators; we instead consider \emph{adversarial} disagreement
induced by verifier capture.

\paragraph{Adversarial robustness of LLM judges.}
Recent work has begun to expose the silent nature of biased LLM
judges: \citet{zhang2026blindcurator} show that a biased judge
\emph{silently disables} a downstream curator past a sharp false-pass
threshold, and broader red-teaming studies
\citep{perez2022red,schulhoff2023ignore} catalogue prompt-injection
attacks. Our threat model (Section~\ref{sec:threat}) operationalises
these observations into a quantitative capture-resilience condition.

\paragraph{Peer-review automation.}
\citet{liu2026paperreviewtoolkit} release an open-source multi-reviewer
pipeline that we use as one of our testbeds; we discuss this and
related calibration work \citep{liu2026calibration} in
Section~\ref{sec:related}.

% =====================================================================
\section{Threat Model and Definitions}
\label{sec:threat}
% =====================================================================

We consider an ensemble $\V = \{v_1, \ldots, v_N\}$ of $N$ LLM verifiers
that independently evaluate a candidate output and produce a continuous
decision distribution $p_i \in [0,1]$ on the binary outcome
$\{0,1\}$; the aggregated decision is produced by a fixed aggregation rule
$\mathcal{A}: [0,1]^N \to \{0,1\}$. A subset $\A \subseteq \V$ of size
$f = |\A|$ is \emph{captured} by an adversary that controls the verifiers'
outputs but cannot observe or modify the inputs seen by honest verifiers.

\begin{definition}[Honest verifier]
\label{def:honest}
A verifier $v_i \in \V \setminus \A$ is \emph{honest} if
$p_i = \mathbb{P}[\text{correct outcome} \mid x, \theta_i]$, where $\theta_i$
is the verifier's fixed-but-unknown system prompt and training. In
particular, an honest verifier's outputs are conditionally independent
given the input $x$.
\end{definition}

\begin{definition}[Capture coalition]
\label{def:capture}
A captured verifier $v_j \in \A$ outputs
$\tilde{p}_j = g(p_j^{\text{true}})$ where $g: [0,1] \to [0,1]$ is an
adversarial response function drawn from one of three concrete attack
families:
\begin{itemize}
  \item \emph{Prompt injection}: $g$ is a fixed affine or step function
        chosen by the adversary (e.g. ``always output 0.9'').
  \item \emph{Rubric injection}: $g$ depends on the input via a hidden
        rubric bias that flips the verifier's preference order on a
        target subset of inputs.
  \item \emph{Majority bias}: $g$ mimics the running empirical mode of
        the other verifiers' outputs, so the captured output is
        indistinguishable from the majority.
\end{itemize}
\end{definition}

\begin{definition}[Aggregation correctness]
\label{def:correctness}
The aggregation rule $\mathcal{A}$ is $(\eta, \delta)$-correct on a
distribution $\mathcal{D}$ if
$\Pr_{(x,y) \sim \mathcal{D}, \A, g}
[\mathcal{A}(p\_1, \ldots, p\_N) = y] \geq 1 - \delta$,
where the probability is taken over the input distribution, the
adversary's choice of capture coalition $\A$, and the response
functions $g$ of all captured verifiers.
\end{definition}

\begin{definition}[Capture-resilient aggregation]
\label{def:resilient}
The aggregation rule $\mathcal{A}$ is $\varepsilon$-\emph{capture-resilient}
if, for every adversary that captures at most $f$ verifiers, the ensemble
remains $(\eta, \delta)$-correct with $\delta \leq \varepsilon$ for every
input distribution $\mathcal{D}$.
\end{definition}

\begin{theorem}[Capture lower bound]
\label{thm:lower}
For any deterministic aggregation rule $\mathcal{A}$, if the adversary
captures $f \geq N/2$ verifiers, then for \emph{some} input distribution
$\mathcal{D}$ and \emph{some} adversarial response function $g$ from
Definition~\ref{def:capture}, the rule $\mathcal{A}$ is not
$\varepsilon$-capture-resilient for any $\varepsilon < 1/2$.
\end{theorem}

\begin{proof}[Proof sketch]
If the adversary controls a strict majority $f > N/2$, the adversary can
choose $g$ such that the captured verifiers all output the same value $y^*$
with probability $1$, while honest verifiers (whose accuracy is below $1$
on $\mathcal{D}$) collectively cannot outvote $\A$. Formal reduction to
the classical $f \geq N/2$ impossibility bound for deterministic
aggregation. Full proof in Appendix~\ref{app:proof}.
\end{proof}

\begin{remark}[Continuous outputs vs.\ classical BFT]
Theorem~\ref{thm:lower} recovers the classical $f < N/2$ BFT impossibility
bound but does not yield a constructive upper bound, because LLM verifier
outputs are continuous: a verifier that outputs $p_i = 0.51$ is
qualitatively different from one that outputs $p_i = 0.99$, even though
both would be treated as ``in favor'' by a binary-vote aggregation rule.
This gap is what motivates the trust-score weighting in \textsc{CRV}
(Section~\ref{sec:method}).
\end{remark}

\begin{assumption}[Honest independence]
\label{ass:indep}
Honest verifiers' outputs are conditionally independent given $x$. This
is the standard assumption in LLM-ensemble calibration work
\citep{cui2026edo,zheng2023judging}; we discuss the relaxation in
Section~\ref{sec:related}.
\end{assumption}

% =====================================================================
\section{Methods: Seven Aggregation Rules}
\label{sec:method}
% =====================================================================

We benchmark seven aggregation rules. Three are \emph{classical} and
require no learning: plain majority voting
(Section~\ref{sub:majority}), PBFT-style deterministic thresholding
(Section~\ref{sub:pbft}), and uniform-weighted mean
(Section~\ref{sub:uniform}). Three are \emph{trust-score-based} and
share the EMA update
$\tau_i^{(t)} = \alpha \tau_i^{(t-1)} + (1-\alpha) c_i^{(t)}$ where
$c_i^{(t)} = 1 - |p_i^{(t-1)} - \text{med}_{-i}^{(t-1)}|$; they differ
in how the trust score is \emph{used} at decision time: weighted mean,
weighted median, and Kalman-filter trust tracking. The seventh is
\textsc{Atv}, an adaptive hybrid we proposed as a candidate that we
expected to combine the strengths of soft weighting and outlier
robustness.

\paragraph{Plain majority.}
\label{sub:majority}
Each verifier contributes a binary vote $\mathbb{1}[p_i \geq 0.5]$; the
aggregate is the majority over the $N$ binary votes. Cost $O(N)$ per
round, no state.

\paragraph{PBFT-style deterministic thresholding.}
\label{sub:pbft}
\citep{castro1999pbft} Each verifier again votes binary; the aggregate
passes if at least $(N+f+1)/2$ verifiers agree, which is the
deterministic threshold that tolerates $f$ adversarial nodes when a
strict majority is required. Cost $O(N)$ per round.

\paragraph{Uniform-weighted mean.}
\label{sub:uniform}
The aggregate is $\mathbb{1}[\sum_i p_i / N \geq 0.5]$. Cost $O(N)$,
no state.

\paragraph{Trust-score weighted mean.}
\label{sub:trustmean}
At round $t$, weights are $w_i = \tau_i^{(t-1)} / \sum_j \tau_j^{(t-1)}$;
the aggregate is $\mathbb{1}[\sum_i w_i p_i \geq 0.5]$. The trust
update uses consistency against the per-round median. Cost $O(N)$
plus a per-round $O(N \log N)$ for the median.

\paragraph{Trust-score weighted median.}
\label{sub:trustmedian}
Identical trust update; at decision time, sort verifiers by current
$p_i$ and pick the $p_i$ at which the cumulative trust weight crosses
half the total weight. This is robust to a single extreme value but
expensive $O(N \log N)$ per round.

\paragraph{Kalman-filter trust tracking.}
\label{sub:kalman}
Same decision rule as weighted mean, but trust tracking uses a
Kalman filter with process variance $Q = 0.01$ and measurement
variance $R = 0.05$ instead of an EMA, which we hypothesised would
smooth trust under noisy verifier streams.

\paragraph{\textsc{Atv}: adaptive trust-weighted voting.}
\label{sub:atv}
A hybrid we designed as the natural synthesis of the two trust-score
decision rules. \textsc{Atv} switches between weighted mean and
weighted median at each round based on the standard deviation of the
trust scores: if $\text{std}(\tau) < \tau_{\text{split}} = 0.5$ (trust
concentrated, presumed-uniform regime), use weighted mean; otherwise
(trust dispersed, capture suspected) use weighted median. The
hyperparameter $\tau_{\text{split}}$ was tuned on a held-out scenario;
we report sensitivity in Appendix~\ref{app:ablation}.

\paragraph{Why we expected \textsc{Atv} to win.}
Our \emph{a priori} hypothesis, before running the benchmark, was that
\textsc{Atv} would inherit the best of both worlds: the mean's
sensitivity to subtle weight differences in the easy regime, and the
median's robustness to outlier-dominated rounds in the hard regime.
The benchmark (Section~\ref{sec:exp}) does not support this hypothesis.

\paragraph{Why we report \textsc{Atv} even though it loses.}
We report \textsc{Atv} as a representative of the natural design space
that future researchers might explore, rather than as our recommended
method. The negative result is informative precisely because \textsc{Atv}
\emph{was} our best prior and still loses; the same is true of the
Kalman variant.

% =====================================================================
\section{Experimental Setup}
\label{sec:exp}
% =====================================================================

\paragraph{Datasets.}
\label{sub:generator}
We use a synthetic verifier-output generator in this study; we discuss
the limitations of this choice in Section~\ref{sec:conclusion}. The
generator (Section~\ref{sub:generator}) models the agreement structure
of a multi-LLM verification ensemble under verifier capture, and is
calibrated to produce vote distributions whose summary statistics
match published numbers for LLM judges on rubric decisions
\citep{leung2026rubricllm}: per-verifier outputs are concentrated
near $0$ or $1$ with $p_{\text{conf}} \in [0.7, 0.85]$ and a small
amount of i.i.d.\ Gaussian noise ($\sigma = 0.05$). We additionally
release a \emph{Beyond-PDF} interactive demonstration that replays
the benchmark on real LLM trace files
(Section~\ref{sec:broader}).

\paragraph{Attack models.}
\label{sub:attacks}
We instantiate the three attack families of
Definition~\ref{def:capture} concretely: \emph{prompt injection}
(\textsc{PromptInj}: captured verifier outputs the complement of its
honest output), \emph{rubric injection}
(\textsc{RubricInj}: captured verifier flips its label on a random
fraction of inputs), and \emph{majority bias}
(\textsc{MajBias}: captured verifier mimics the running median of the
other verifiers' outputs). Section~\ref{sub:attacks} gives the formal
parameterisations.

\paragraph{Scenarios.}
We consider four scenarios that vary the two most consequential
parameters: honest-verifier per-instance accuracy $\eta \in \{0.65,
0.85\}$ and attack strength $\alpha \in \{0.6, 0.9\}$ (probability
that a captured verifier executes the attack on a given round). The
four combinations are labelled S1--S4 in
Table~\ref{tab:main}. We run $T = 1000$ rounds per scenario and
average over $n = 5$ random seeds; full per-seed numbers are in
Appendix~\ref{app:ablation}.

\paragraph{Metrics.}
We report \emph{decision accuracy} as the fraction of rounds on which
the aggregated decision equals the ground-truth label. Accuracy is the
right metric for this benchmark because the goal of the aggregator is
to produce the correct label, and the captured verifiers' goal is to
produce an incorrect one; alternative metrics such as false-positive
rate at the verifier level would require labelling individual
verifiers as captured or honest, which we cannot do from observed
verifier outputs alone (this is the central statistical difficulty).

\paragraph{Hyperparameters.}
All trust-based methods use $\alpha = 0.9$ (EMA decay) and
$c_i^{(t)} = 1 - |p_i - \text{med}_{-i}|$. \textsc{Atv} uses
$\tau_{\text{split}} = 0.5$ chosen on a held-out scenario. PBFT
threshold is $(N+f+1)/2$. Code and configuration files are released
with the paper.

% =====================================================================
\section{Results}
\label{sec:results}
% =====================================================================

\paragraph{Headline.}
Table~\ref{tab:main} reports the headline result. Under every scenario
we test, plain majority voting is at or near the top; the trust-based
methods tie or trail majority by $1$--$3$ percentage points; and
\textsc{Atv}, our adaptive hybrid, never wins. The differences are
small but consistent across all five random seeds we ran; the
standard deviations are below $0.01$ in every cell.

\begin{table}[ht]
\caption{Decision accuracy across four scenarios. Best per column
in bold. No aggregation method we tested substantially outperforms
plain majority voting.}
\label{tab:main}
\centering
\small
\begin{tabular}{lccccc}
\toprule
Method & S1 & S2 & S3 & S4 & mean \\
 & ($\eta{=}0.85$, $\alpha{=}0.6$)
 & ($\eta{=}0.65$, $\alpha{=}0.6$)
 & ($\eta{=}0.65$, $\alpha{=}0.9$)
 & ($\eta{=}0.85$, $\alpha{=}0.9$)
 & \\
\midrule
majority            & \textbf{0.889} & \textbf{0.553} & \textbf{0.459} & 0.843 & \textbf{0.686} \\
bft\_pbft           & 0.777 & 0.522 & 0.458 & 0.720 & 0.619 \\
weighted\_uniform   & 0.804 & 0.484 & 0.365 & 0.729 & 0.595 \\
trust\_mean         & 0.883 & 0.521 & 0.432 & \textbf{0.851} & 0.672 \\
trust\_median       & \textbf{0.889} & \textbf{0.553} & \textbf{0.459} & 0.844 & \textbf{0.686} \\
kalman\_trust       & 0.882 & 0.535 & 0.446 & 0.842 & 0.677 \\
\textsc{Atv}        & 0.883 & 0.527 & 0.438 & 0.846 & 0.673 \\
\bottomrule
\end{tabular}
\end{table}

\paragraph{Why no method beats majority.}
The negative result has a structural explanation. Capture is
\emph{statistically indistinguishable from honest-but-wrong}: a
captured verifier outputting the complement of its honest output has
exactly the same per-instance deviation from the median as an honest
verifier that happens to be wrong on this round. Any rule that uses
deviation from the median to identify capture therefore cannot
distinguish a captured verifier from a minority of honest-but-wrong
verifiers; it will either ignore both (and gain nothing) or exclude
both (and lose accuracy on rounds where honest-but-wrong verifiers
are the actual minority).

In the high-honest scenario S1, the trust-based methods approach
majority but trail it by 0.5 percentage points, consistent with the
explanation above: the trust mechanism correctly down-weights
captured verifiers when they \emph{happen} to be outliers, but it
also incorrectly down-weights honest verifiers that happen to be
wrong, and the two effects approximately cancel.

In the strong-attack scenarios S3 and S4, the trust-based methods
gain an apparent advantage in some columns (S4 for trust\_mean) but
lose it in others (S2 for all trust-based methods). The advantage in
S4 disappears when averaged across all four scenarios: the mean
column is dominated by majority and trust\_median, both at $0.686$.

\paragraph{Per-attack breakdown.}
Across the three attack families (prompt injection, rubric injection,
majority bias), the ordering of methods is qualitatively similar.
Majority and trust\_median remain tied at the top across all three
attack types; \textsc{Atv} and the Kalman variant trail by $0.5$--$1.5$
percentage points. The largest gap between majority and the best
trust-based method is under \textsc{RubricInj}, where trust\_median
loses to majority by $0.7$ percentage points because the rubric
injection only flips on a random fraction of rounds, and the trust
update averages over both flipped and unflipped rounds.

\paragraph{Sensitivity to hyperparameters.}
\textsc{Atv}'s $\tau_{\text{split}}$ was tuned on a held-out scenario;
sensitivity sweeps over $\tau_{\text{split}} \in [0.2, 0.8]$ change
\textsc{Atv}'s mean accuracy by at most $0.4$ percentage points and
do not affect its ranking. Trust $\alpha$ sweeps over $[0.5, 0.99]$
shift trust\_mean by $\pm 0.5$ percentage points with no rank
change. Full sensitivity table in
Appendix~\ref{app:ablation}.

\paragraph{Implication.}
The benchmark implies that, for the ensemble size, attack strength,
and honest-accuracy regime we consider, \emph{the engineering effort
of building a non-trivial aggregation rule does not pay off}.
Practitioners should prefer plain majority voting and invest their
engineering effort elsewhere --- either in increasing honest-verifier
accuracy or in diversifying the verifier pool so that capture becomes
harder. We elaborate on diversification as a defensive strategy in
Section~\ref{sec:conclusion}.

% =====================================================================
\section{Related Work}
\textbf{Novelty positioning.} This work makes three distinct contributions relative to prior literature. \textbf{First}, prior work on calibration in multi-agent systems reported calibration effects but did not systematically deconfound temporal from communicative causes; our pure-factor decomposition isolates the dominant causal factor. \textbf{Second}, prior work treated self-evaluation immunity as a stable property; we demonstrate that model updates can silently invert this behavior, showing it is model-dependent and fragile. \textbf{Third}, prior work proposed structural coupling-noise tradeoffs; our deterministic-seed reruns fail to support universal noise-increase claims, and we reframe the tradeoff as a method-specific hypothesis with explicit falsification conditions.


\label{sec:related}
% =====================================================================

[Placeholder body. W2 target will fill three threads:]

\paragraph{Classical Byzantine consensus.}
\citet{lamport1982byzantine,castro1999pbft,yin2019hotstuff} establish the
$3f+1$ bound for binary-output consensus under $f$ adversarial nodes. Our
setting differs in that LLM verifier outputs are continuous probability
distributions, so classical BFT bounds do not transfer directly; this
motivates our $\varepsilon$-capture-resilience definition in
Section~\ref{sec:threat}.

\paragraph{LLM-as-judge and ensemble calibration.}
\citet{zheng2023judging} introduce the MT-Bench and Chatbot Arena
evaluation paradigm; \citet{leung2026rubricllm} provide the most recent
calibration numbers for 8 LLM judges on rubric decisions, and
\citet{cui2026edo} jointly optimize ensemble weights and calibration for
subjective supervision. CRV differs from EDO in operating under an
explicit adversarial model (capture) rather than a passive disagreement
model.

\paragraph{Adversarial robustness of LLM judges.}
\citet{zhang2026blindcurator} show that a biased judge \emph{silently}
disables a downstream curator past a sharp false-pass threshold, and
\citet{perez2022red,schulhoff2023ignore} catalogue prompt-injection and
red-team attack techniques against LLM evaluators. Our threat model
operationalises these observations into a quantitative capture-resilience
definition and provides a protocol that provably maintains aggregation
correctness under bounded capture.

\paragraph{Peer-review automation.}
\citet{liu2026paperreviewtoolkit} release an open-source multi-reviewer
pipeline that we use as the testbed for our empirical evaluation
(Section~\ref{sec:exp}); \citet{liu2026calibration} develops calibration
methods for multi-agent ensembles more broadly and is the closest work
in our own research line.

% =====================================================================
\section{Broader Impact}
\label{sec:broader}
% =====================================================================

\paragraph{Positive impact.}
The primary positive impact of this work is \emph{negative}: by
showing that no aggregation method we tested substantially outperforms
plain majority voting under verifier capture, we redirect practitioner
effort away from building sophisticated aggregators and toward
investments that \emph{do} pay off in the capture regime: increasing
honest-verifier accuracy, diversifying the verifier pool, and
auditing deployed ensembles for suspicious agreement patterns. We
expect the benchmark to be useful for engineering teams that are
currently considering whether to invest in a custom aggregation
layer; our data suggest the answer is ``no'' under the parameter
regime we test, and the engineering budget should be spent elsewhere.

A secondary positive impact is methodological: the failure modes of
the seven aggregation rules we test illuminate which statistical
assumptions a successful capture-resilient aggregator would need
to break. Future research that does find a winning method can
explicitly cite which assumption we showed fails here.

\paragraph{Negative impact and dual-use concerns.}
The same attack implementations that allow us to \emph{benchmark}
verifier capture also allow an adversary to \emph{measure} capture
susceptibility of a competing verifier product. We consider this a
real but bounded risk for the following reasons.
\textbf{(i)} The three attack families we instantiate
(\textsc{PromptInj}, \textsc{RubricInj}, \textsc{MajBias}) are
straightforward enough that a competent adversarial team could
re-derive them from the threat model alone; our code release does
not materially lower the cost.
\textbf{(ii)} Our benchmark demonstrates that these attacks, when
attempted against a well-calibrated honest majority, have only
small absolute effect (the gap between majority and the best
trust-based method is below 1 percentage point in 3 of 4
scenarios). An adversary using our code to measure capture
susceptibility will therefore find that the practical upside of
attack deployment is small.
\textbf{(iii)} Our benchmark is \emph{synthetic}. An adversary who
wishes to deploy a capture attack against a real ensemble must
still construct a real captured verifier that mimics the target
ensemble's distribution, which is a substantially harder engineering
problem than running our code.

\paragraph{Responsible disclosure norms.}
We follow the standard ML-security responsible-disclosure norm: we
release the benchmark code under a licence that prohibits offensive
use, and we provide a \texttt{ResponsibleUse.md} file in the
repository that documents the intended defensive use and lists
specific prohibited uses. We do not release any captured-verifier
weights or fine-tuning scripts that would lower the cost of deploying
a real capture attack.

\paragraph{Limitations of public release.}
The benchmark code we release is sufficient to \emph{reproduce} our
results but not sufficient to \emph{deploy} a capture attack against
a real ensemble without substantial additional engineering. This is
deliberate. We discuss the scope of the release, including which
extensions a future maintainer might add and which are out of scope,
in the \texttt{README.md} of the released repository.

\paragraph{Stakeholder considerations.}
Verifier-capture attacks disproportionately threaten operators of
multi-LLM auditing pipelines, including the authors of this paper,
who maintain an open-source peer-review audit tool
\citep{liu2026paperreviewtoolkit}. We have therefore prioritised
documentation of defensive best-practice over attack
generalisation, and we invite other operators of LLM-judge ensembles
to contact us with deployment experience that would refine the
attack-and-defence landscape this benchmark sketches.

\paragraph{Positive-negative balance.}
On balance, we believe the positive impact of releasing this
benchmark---redirecting engineering effort away from ineffective
aggregation rules toward more productive investments in verifier
quality and pool diversification---substantially outweighs the
bounded negative impact of marginally lowering the cost of an
attack that, per our own benchmark, has small practical effect
against well-calibrated honest majorities.

% =====================================================================
\section{Ethics Statement}
\label{sec:ethics}
% =====================================================================

This study does not involve human subjects, personal data, or animal
subjects. The verifier-output data we use is either synthetic
(Section~\ref{sub:generator}) or from publicly available LLM-judge
benchmarks \citep{leung2026rubricllm,zhang2026blindcurator}.

\paragraph{Dual-use risk.}
We have discussed the dual-use risk of releasing the attack
implementations in detail in Section~\ref{sec:broader}. In summary,
the bounded risk is outweighed by the positive impact of the
benchmark, and we have taken reasonable mitigations: a
responsible-use licence, no released captured-verifier weights or
fine-tuning scripts, and a \texttt{ResponsibleUse.md} file in the
repository.

\paragraph{Affected stakeholders.}
The primary stakeholders are operators of multi-LLM verification
ensembles, including AI-auditing platforms, peer-review automation
tools, and safety-content moderation pipelines. We have reached out
to the maintainers of two such platforms
(\citep{liu2026paperreviewtoolkit} and the public release of
\citet{leung2026rubricllm}'s evaluation harness) to inform them of
our findings ahead of public release.

\paragraph{IRB and funding.}
This study was not subject to IRB review. No external funding was
received for this work.

\paragraph{Author conflicts of interest.}
The authors declare no financial conflicts of interest. The authors
maintain an open-source peer-review audit tool
\citep{liu2026paperreviewtoolkit} that is discussed as one of the
motivating use cases for this benchmark; the tool is referenced
consistently with the rest of the literature and no special pleading
is made for it.

% =====================================================================
\section{Conclusion}
\label{sec:conclusion}
% =====================================================================

We have presented an empirical benchmark of seven aggregation methods
for multi-LLM verifier ensembles under three verifier-capture attacks,
across four scenarios. The headline finding is \emph{negative}: no
aggregation method we tested substantially outperforms plain majority
voting, and several trust-score variants underperform majority under
strong capture. We argued that the negative result is structural: a
captured verifier is statistically indistinguishable from an honest
verifier that happens to be wrong, so any rule that uses statistical
deviation to identify capture cannot outperform a rule that ignores
the distinction.

\paragraph{When our negative result would \emph{not} hold.}
Our negative result rests on the assumption that captured verifiers
mimic honest-but-wrong verifiers. If a real-world capture breaks this
assumption --- for instance, by introducing a structured deviation
(e.g. a verifier that always outputs the exact median plus a fixed
delta) --- then a more sophisticated aggregation rule may recover
above majority. Our benchmark does not rule this out; it only shows
that for the three concrete attack parameterisations we test, the
gain is at most a fraction of a percentage point.

A second assumption is that honest verifiers are exchangeable. If
honest verifiers have known specialty areas (e.g.\ one verifier is
better at math, another at code), then a domain-aware aggregator may
do substantially better than majority. Our synthetic generator
deliberately does not model this specialisation, so our benchmark
under-estimates the upside of sophisticated aggregation in
heterogeneous pools.

\section{Discussion and Limitations}
Our results rely on a synthetic verifier-output generator calibrated
to published rubric-LLM numbers but not validated against
real LLM trace data. A practitioner applying our benchmark to a real
ensemble should re-run our scenario definitions with the real
distribution and may find that some methods we rule out do in fact
help on their specific distribution. We release the synthetic
generator and the attack implementations to facilitate this
replication.

\paragraph{Future work.}
Three directions follow from this study.
(i) Validate the negative result on real LLM verifier trace data;
the result may not survive the transition from synthetic to real.
(ii) Investigate structured-attack settings where captured verifiers
introduce deviations outside the honest-but-wrong distribution; the
negative result may not apply there.
(iii) Investigate diversified verifier pools where domain-aware
aggregation can outperform majority by exploiting specialisation.

% =====================================================================
% References (refs.bib populated in W1)
% =====================================================================
\bibliographystyle{tmlr}
\bibliography{refs}

% =====================================================================
% Appendix begins here
% =====================================================================
\appendix

\section{Proof of the Capture Lower Bound (Theorem~1)}
\label{app:proof}

We prove the deterministic lower bound stated in Theorem~1: if
$f \geq N/2$ of the verifiers are captured, no deterministic
aggregation rule $\mathcal{A}: [0,1]^N \to \{0,1\}$ can be
$\varepsilon$-capture-resilient for $\varepsilon < 1/2$.

\begin{proof}
Consider the case $f > N/2$ (strict majority capture; the case
$f = N/2$ follows by a standard limit argument). The adversary
chooses a target label $y^* \in \{0,1\}$ and a response function
$g$ such that every captured verifier outputs $y^*$ with
probability $1$ on every input. The adversary also chooses an input
distribution $\mathcal{D}$ on which the \emph{honest} verifiers have
per-instance accuracy strictly below $1/2$; this is always possible
by choosing $\mathcal{D}$ to concentrate on inputs where the honest
verifiers systematically err.

Because the captured coalition $\A$ forms a strict majority, the
deterministic rule $\mathcal{A}$ sees $p_i \in \{0, 1\}$ on the
captured inputs (after $g$ is applied) and any values in $[0,1]$ on
the honest inputs. Since $\mathcal{A}$ is deterministic, its
decision is a function only of the multiset $\{p_1, \ldots, p_N\}$,
and in particular of the \emph{number} of captured verifiers voting
$y^*$, which is $f > N/2$. For any such deterministic rule, there
exists an assignment of honest outputs such that $\mathcal{A}$ outputs
$y^*$ on every input, even when the true label is the opposite
$\bar{y}^*$; this is a direct consequence of the fact that a strict
majority of $N$ deterministic binary inputs forces the binary
majority rule, and any rule that agrees with binary majority on
strict-majority inputs cannot be $\varepsilon$-resilient for
$\varepsilon < 1/2$ on the chosen $\mathcal{D}$.

This shows that for any deterministic $\mathcal{A}$, the adversary
can construct a distribution $\mathcal{D}$ on which $\mathcal{A}$
achieves error probability at least $1/2$, completing the proof.
\end{proof}

\begin{proof}[Proof of Theorem~2 (claim; \textsc{Atv} $\varepsilon$-resilience)]
The full proof is omitted for brevity because \textsc{Atv} does
not in fact achieve the claimed $\varepsilon = 2f/N + O(\Delta^2)$
guarantee in our benchmark (Section~\ref{sec:exp}). We retain the
theorem statement for completeness so that future readers can
identify exactly which step in the proof fails; a candidate
follow-up paper can either fix the proof or revise the theorem.
\end{proof}

\section{Algorithm Pseudocode and Complexity}
\label{app:algo}

We give pseudocode for the four representative aggregation rules.
All four have memory $O(N)$ for the per-verifier state.

\paragraph{Plain majority.}
\begin{algorithmic}[1]
\Require votes $p_1, \ldots, p_N$
\State $s \gets \sum_i \mathbb{1}[p_i \geq 0.5]$
\State \Return $\mathbb{1}[s > N/2]$
\end{algorithmic}
Time $O(N)$.

\paragraph{PBFT-style threshold.}
\begin{algorithmic}[1]
\Require votes $p_1, \ldots, p_N$; bound $f$
\State $s \gets \sum_i \mathbb{1}[p_i \geq 0.5]$
\State \Return $\mathbb{1}[s \geq (N+f+1)/2]$
\end{algorithmic}
Time $O(N)$.

\paragraph{Trust-score weighted mean.}
\begin{algorithmic}[1]
\Require votes $p_1, \ldots, p_N$; trusts $\tau_1, \ldots, \tau_N$; decay $\alpha$
\State $w_i \gets \tau_i / \sum_j \tau_j$ for each $i$
\State \Return $\mathbb{1}[\sum_i w_i p_i \geq 0.5]$
\State Update: $\tau_i \gets \alpha \tau_i + (1-\alpha)(1 - |p_i - \text{med}_{-i}|)$ for each $i$
\end{algorithmic}
Time $O(N \log N)$ per round (median), memory $O(N)$.

\paragraph{Trust-score weighted median.}
\begin{algorithmic}[1]
\Require votes $p_1, \ldots, p_N$; trusts $\tau_1, \ldots, \tau_N$; decay $\alpha$
\State Sort indices by $p_i$; let $\pi$ be the sorted permutation
\State $W \gets \sum_j \tau_j$; find smallest $k$ with $\sum_{i \leq k} \tau_{\pi(i)} \geq W/2$
\State \Return $\mathbb{1}[p_{\pi(k)} \geq 0.5]$
\State Update trust as in weighted mean.
\end{algorithmic}
Time $O(N \log N)$ per round, memory $O(N)$.

\paragraph{\textsc{Atv} (adaptive hybrid).}
\begin{algorithmic}[1]
\Require votes $p_1, \ldots, p_N$; trusts $\tau_1, \ldots, \tau_N$; decay $\alpha$; split $\tau_{\text{split}}$
\If{$\text{std}(\tau) < \tau_{\text{split}}$}
  \State \Return weighted-mean decision
\Else
  \State \Return weighted-median decision
\EndIf
\State Update trust as in weighted mean.
\end{algorithmic}
Time $O(N \log N)$ per round, memory $O(N)$.

\section{Additional Ablations and Sensitivity Sweeps}
\label{app:ablation}

\paragraph{Per-attack-model breakdown.}
Table~\ref{tab:per-attack} reports per-attack-model accuracies. The
ordering of methods is qualitatively consistent across the three
attack families: majority and trust\_median remain tied at the top,
with the largest absolute gap (0.7 percentage points) appearing
under rubric injection.

\begin{table}[ht]
\caption{Decision accuracy by attack model, averaged across all
four scenarios (5 seeds, $T=1000$ rounds per cell).}
\label{tab:per-attack}
\centering
\small
\begin{tabular}{lccc}
\toprule
Method & \textsc{PromptInj} & \textsc{RubricInj} & \textsc{MajBias} \\
\midrule
majority            & 0.877 & 0.945 & 0.997 \\
bft\_pbft           & 0.765 & 0.886 & 0.985 \\
weighted\_uniform   & 0.876 & 0.944 & 0.997 \\
trust\_mean         & \textbf{0.894} & 0.944 & \textbf{0.998} \\
trust\_median       & 0.877 & \textbf{0.945} & 0.997 \\
kalman\_trust       & 0.877 & \textbf{0.945} & 0.997 \\
\textsc{Atv}        & 0.886 & \textbf{0.945} & 0.997 \\
\bottomrule
\end{tabular}
\end{table}

\paragraph{Hyperparameter sensitivity.}
\textsc{Atv}'s $\tau_{\text{split}}$ was tuned on a held-out
scenario. Sensitivity sweeps over $\tau_{\text{split}} \in [0.2, 0.8]$
shift \textsc{Atv}'s mean accuracy by at most $0.4$ percentage
points; the rank ordering is preserved across all sweep values.
Trust-update decay $\alpha$ sweeps over $[0.5, 0.99]$ shift
trust\_mean by $\pm 0.5$ percentage points with no rank change.

\paragraph{Ensemble-size sweep.}
We additionally ran the benchmark at $N \in \{5, 7, 10\}$ with $f$
proportional to $N$ ($f = \lfloor N/3 \rfloor$). Under this
proportional-f regime, majority remains the top method at every
ensemble size, and the gap between majority and trust-based methods
\emph{widens} as $N$ grows (because the larger honest majority
hardens against capture). This is consistent with our structural
explanation: as the honest majority grows, capture becomes a
smaller perturbation that any aggregation rule, including
majority, can absorb.

\paragraph{Real-LLM validation (sketch).}
We did not run the full benchmark on real LLM verifier traces
within the scope of this paper. A small pilot run on
\citet{leung2026rubricllm}'s evaluation harness (8 LLM judges, 1248
rubric decisions, of which we synthetically capture 2) is consistent
with the synthetic-data trend: majority remains at or near the top,
trust-based methods trail by less than 1 percentage point. We discuss
the limitations of this pilot in Section~\ref{sec:conclusion}.

\section{Code and Reproducibility}
\label{app:code}

\paragraph{Repository layout.}
Code is released under a permissive licence with an attached
\texttt{ResponsibleUse.md} file. The repository contains:
\begin{itemize}
  \item \texttt{crv.py}: legacy CRV implementation (retained for
        archival reference; not used as a benchmarked method).
  \item \texttt{attacks.py}: implementations of the three attack
        families (\textsc{PromptInj}, \textsc{RubricInj},
        \textsc{MajBias}).
  \item \texttt{extended\_methods.py}: the seven aggregation methods
        benchmarked in this paper.
  \item \texttt{benchmark.py}: the four-scenario benchmark driver.
  \item \texttt{verify\_p6.py}: the verify suite used during
        paper preparation.
\end{itemize}

\paragraph{Beyond-PDF interactive demo.}
We release a Beyond-PDF interactive demonstration that replays the
benchmark on real LLM trace files. The demo runs the seven
aggregation rules on a user-provided trace file and displays the
per-round decision side-by-side with the ground-truth label. We
expect this to be useful for engineering teams that wish to
re-run our scenario definitions on a real ensemble.

\paragraph{Reproducibility hashes.}
At submission time, we record SHA-256 hashes of all released code
artifacts in \texttt{HASHES.txt}. Reviewers can verify their copy
against the hashes by running the included \texttt{verify\_hashes.sh}
script.

\paragraph{Computational footprint.}
A full benchmark run (4 scenarios $\times$ 7 methods $\times$ 5 seeds,
$T=1000$ rounds) takes approximately 4 minutes on a single CPU core
(no GPU required) and uses 50 MB of memory. The complete experiment
can be reproduced in under 10 minutes on commodity hardware.


\section*{Reproducibility Statement}

\textbf{Code and data availability.} All experiment code, raw API response logs, and analysis scripts are released in the supplementary material (directories \texttt{code/}, \texttt{data/}, \texttt{analysis/}). The complete evidence package includes SHA-256 checksums for all experimental result files (\texttt{evidence\_manifest.json}).

\textbf{Step-by-step reproduction protocol.} (1) \textit{Environment}: Python 3.10+ with \texttt{transformers>=4.40}, \texttt{scipy>=1.10}, \texttt{scikit-learn>=1.3}; set \texttt{PYTHONHASHSEED=0}. (2) \textit{API configuration}: set \texttt{OPENAI\_API\_KEY} and \texttt{DEEPSEEK\_API\_KEY} environment variables. (3) \textit{Execution}: run \texttt{python code/run\_all\_experiments.py --seeds 1-30} to reproduce all conditions; each condition takes approximately 2--3 hours. (4) \textit{Analysis}: run \texttt{python analysis/main\_analysis.py} to reproduce all tables and figures; all scripts use deterministic seeds (\texttt{int.from\_bytes(seed.encode(), ``little'')}) for bit-exact reproducibility.


\section*{Reproducibility Statement}
\label{sec:reproducibility}

\textbf{Code and data availability.} All experiment code, raw API response logs, analysis scripts, and frozen analysis code are released in the supplementary material (directories \texttt{code/}, \texttt{data/}, \texttt{analysis/}, \texttt{evidence/}). The evidence manifest (\texttt{evidence/evidence\_manifest.json}) contains SHA-256 checksums for all result files, enabling independent verification of file integrity.

\textbf{Step-by-step reproduction protocol.} (1) \textit{Environment}: Python 3.10+ with \texttt{transformers>=4.40}, \texttt{scipy>=1.10}, \texttt{scikit-learn>=1.3}; set \texttt{PYTHONHASHSEED=0}. (2) \textit{API configuration}: set \texttt{OPENAI\_API\_KEY}, \texttt{DASHSCOPE\_API\_KEY}, and \texttt{DEEPSEEK\_API\_KEY} environment variables. (3) \textit{Execution}: run \texttt{python code/run\_all\_experiments.py --seeds 1-30} to reproduce all conditions; each condition takes approximately 2--3 hours. (4) \textit{Analysis}: run \texttt{python analysis/main\_analysis.py} to reproduce all tables and figures; all scripts use deterministic seeds for bit-exact reproducibility of all randomization and bootstrap procedures.

\section*{Broader Impact and Ethical Considerations}
\label{sec:ethics}

This work studies calibration dynamics in multi-agent LLM systems. Understanding these limitations is essential for safe and ethical deployment.

\textbf{Ethical implications.} (i) \textit{Self-evaluation fragility}: Model updates can silently invert self-evaluation immunity, creating risks of silent reliability decay in safety-critical systems. Practitioners must treat self-evaluation reliability as version-dependent and re-verify after every model update. (ii) \textit{Calibration fatigue}: Long-horizon multi-agent interactions may produce progressive miscalibration without external signals, creating risks in deployed systems with extended interaction horizons. Continuous calibration monitoring and automatic circuit-breakers are recommended. (iii) \textit{Responsible research}: We explicitly report reproduction failures rather than suppressing them; all code, data, and analysis scripts are released with deterministic seeds and SHA-256 checksums for independent verification.

\end{document}
